\begin{abstract}
When a language model writes code for a multi-step task, what does it represent about the operation it is about to perform? We find that information about the next subgoal is available before the model emits its implementation. We localize a small set of attention heads using cross-language activation patching and introduce \goalscope, which uses their output contributions to elicit a natural-language description of the upcoming subgoal in a separate model run. The same contributions can steer a new code context to implement the source subgoal in another programming language. Small head sets support effective readout, while a broader set is causally important for native implementation: preserving its access to task information largely preserves performance, and restricting its attention sharply reduces performance relative to random-head controls. Query and key/value interventions further show that these heads support goal-specific selection and adaptation to changed task information. An OV-based readout traces the availability of subgoal information through model depth. We replicate readout and implementation-transfer experiments within multiple models and explore readouts beyond explicitly specified code tasks. Together, these findings identify a readable, transferable, and causally relevant representation of an upcoming subgoal, without presupposing a particular algorithm for selecting it.
\end{abstract}
